\section{Limitations}\label{sec:limits}

\subsection{Threats to validity}\label{sec:threats}

\paragraph{Internal validity.} The intervals come from \nSeeds{} seeds for the main experiment, \nSeedsAbl{} for the ablation and three for the timing, and a Student interval with so few seeds is wide; a difference whose interval contains zero is not distinguished from none, which is the case for the steps per weight, the weights above 20 and the converged start. The seeds fix the random superposition that starts every method, so the pairing by seed removes the variation of the start,
and what it leaves is the variation of the optimization. The conjugate-gradient iteration is stopped after 15 steps per weight and the loss is not convex, so the solutions are local and the monotonicity of Proposition~\ref{prop:mono}(a), proved for global minimizers, is observed and not guaranteed. The diagnostic of Section~\ref{sec:res-rivals} shows the consequence: the solutions at \diagGswHundredMin{} to \diagGswHundredMax{} of the weights are worse on the loss than a hologram of weighted GS.
The loop of the ablation is a copy of the module's loop; it is checked against the module on the first seed (largest absolute difference \consistMax{}) and for the other seeds it is the same code path, so a difference between the two on another seed is possible but was not looked for. The seconds of the main experiment come from processes with one thread, and the seconds of the timing from a machine shared with other jobs.
The dominance criterion of Eq.~\ref{eq:ndom} is the one of the module's test, with a strict inequality in the uniformity only, so a front point that ties a published point in the efficiency and exceeds it in the uniformity counts, and one that ties in the uniformity does not.

\paragraph{External validity.} The experiment is one task: a ten by ten lattice of equal traps in one plane, on a 256 pixel SLM padded by 2, with one wavelength and one lens, on the ideal device. The benchmark extends the check to five other geometries, but only through the stored fronts and the cell medians, which were not produced for this comparison: the stored front is the front of the first seed with a converged budget that the wall-clock bound shortens, and the published points are medians over a number of seeds that
the same bound sets. The result on the benchmark is mixed: at 100 iterations the stored front meets the union of the published points on \benchPassHundred{} of \benchNHundred{} geometries, so the claim of the paper does not extend beyond the acceptance task. The bake lists four cases with three seeds each. No hardware was used: the uniformities of the paper are those of the scalar model and not of an instrument, and the published points are computed by the engine's implementation of each method, not by the code of its authors.
Arrays with unequal weights, traps in several planes, vortex traps and the device with crosstalk, quantization and flicker are outside what was measured for the method; the weights $w_m$ enter the loss and the metrics, but the bounds of Proposition~\ref{prop:bound} are proved only for equal weights.

\paragraph{Construct validity.} The method optimizes the normalized squared error $E$ and the normalized light $S$, and the paper evaluates the front in the plane of the efficiency $e$ and the uniformity $u$. The uniformity depends on the brightest and dimmest traps only and the error on all traps, so the two are related by the inequalities of Proposition~\ref{prop:bound} and not by a function, and a front of $(E,-S)$ is not guaranteed to be a front in $(e,u)$; the readout keeps the non-dominated points of the second plane, and
the number that remains is \nParetoMean{} of \nLambda{} on average. The hypervolume with the reference $(0,0)$ is dominated by the uniform end; three references were tested and the ordering did not change, but a reference chosen by an experimenter for a particular array can weigh the parts of the front differently. The efficiency is the light in the traps as a fraction of the input power, which is the quantity of the foundational comparison and not the diffraction efficiency of an instrument.

\subsection{Scope and what is not claimed}\label{sec:scope}

The paper claims the following and no more. On the ten by ten task of the engine, on the ideal device, a continuation of the normalized weighted loss from the output of phase-fixed WGS, with 20 weights and 15 conjugate-gradient steps per weight, gives a front whose hypervolume exceeds that of the union of the points of four published methods in every one of \nSeeds{} seeds, and that dominates two of those points in every seed. It does not claim that the front is the Pareto front of the problem: part of it is
dominated by iterates of weighted GS. It does not claim an improvement at the uniform end, over the weighted GS point or over the phase-fixed point, where it dominates them in a minority of seeds. It does not claim an improvement on the geometries of the benchmark, where the stored front meets the published union on \benchPassHundred{} of \benchNHundred{} geometries at 100 iterations. It does not claim a time advantage: the time is of the order of one minute on the machine used and depends on its load. It does not claim a result for a device model, for unequal weights, for traps off the
focal plane or for hardware. The identity of Proposition~\ref{prop:shift} and the bounds of Proposition~\ref{prop:bound} are proved and checked numerically, and they hold for every hologram; the monotonicity of Proposition~\ref{prop:mono}(a) is proved for global minimizers and observed for the local ones; Conjecture~\ref{conj:branch} is supported by the counts of Section~\ref{sec:res-front} and by the cold start of the ablation, and is not proved.
